% GENERATED by scripts/build-lecture-slides.py; do not edit. Portable profile: XeTeX, theme and fonts from the core.
\makeatletter\IfFileExists{core/beamer/beamerthememx.sty}{\def\mxcorepath{core/}\def\input@path{{core/beamer/}{core/}}}{\def\mxcorepath{../core/}\def\input@path{{../core/beamer/}{../core/}}}\makeatother
\documentclass[aspectratio=169,10pt,t]{beamer}
\usetheme{mx}
\usepackage{graphicx,booktabs}
\usepackage[french]{babel}
\setlength{\parskip}{.45em}
\hypersetup{pdfauthor={},pdfsubject={Lecture slides},pdfkeywords={},pdftitle={Organisation des données et statistique}}
\providecommand{\N}{\mathbb{N}}
\providecommand{\Z}{\mathbb{Z}}
\providecommand{\D}{\mathbb{D}}
\providecommand{\Q}{\mathbb{Q}}
\providecommand{\R}{\mathbb{R}}
\providecommand{\C}{\mathbb{C}}
\providecommand{\vect}[1]{\overrightarrow{#1}}
\providecommand{\vi}{\vec{\imath}}
\providecommand{\vj}{\vec{\jmath}}
\providecommand{\vk}{\vec{k}}
\providecommand{\norm}[1]{\left\lVert #1\right\rVert}
\providecommand{\vc}[2]{\begin{pmatrix}#1\\#2\end{pmatrix}}
\providecommand{\deter}[2]{\det\left(#1,#2\right)}
\providecommand{\abs}[1]{\left\lvert #1\right\rvert}
\providecommand{\dd}{\,\mathrm{d}}
\providecommand{\e}{\mathrm{e}}
\providecommand{\Card}{\operatorname{Card}}
\providecommand{\E}{\operatorname{E}}
\providecommand{\V}{\operatorname{V}}
\providecommand{\degre}{^{\circ}}
\providecommand{\un}[1]{\,\mathrm{#1}}
\providecommand{\sci}[2]{#1\times10^{#2}}
\providecommand{\pgcd}{\operatorname{PGCD}}
\providecommand{\ppcm}{\operatorname{PPCM}}
\title{Organisation des données et statistique}
\subtitle{Cours et exemples}
\date{}
\author{}
\begin{document}
\mxtitleframe
\begin{frame}{Objectifs}
Calculer proportions et évolutions, résumer une série par ses indicateurs et interpréter une boîte à moustaches.
\end{frame}
\begin{frame}{Proportion et évolution}
Une proportion est un rapport partie/ensemble, avec effectif total non nul. Le taux d’évolution d’une valeur initiale \(V_i>0\) à une valeur finale \(V_f\) est \(t=(V_f-V_i)/V_i\). Le coefficient multiplicateur est \(1+t\); un taux de \(p\%\) vaut \(p/100\).
\end{frame}
\begin{frame}{Proportion et évolution}
Pour des valeurs strictement positives, deux évolutions successives de taux \(t_1,t_2>-1\) ont pour coefficient \((1+t_1)(1+t_2)\); leur taux global est ce produit moins \(1\). Si \(t>-1\), le taux réciproque est \(1/(1+t)-1\).
\end{frame}
\begin{frame}{Proportion et évolution}
La valeur finale après deux évolutions vaut \(V_i(1+t_1)(1+t_2)\). Pour revenir en arrière, multiplier par l’inverse du coefficient.
\end{frame}
\begin{frame}{Proportion et évolution}
Hausse de \(20\%\) puis baisse de \(20\%\) : coefficient \(1{,}2\times0{,}8=0{,}96\), baisse globale \(4\%\). Passer de \(30\%\) à \(40\%\) est une hausse de \(10\) points de pourcentage, mais une hausse relative de \(1/3\), soit environ \(33{,}3\%\).
\end{frame}
\begin{frame}{Indicateurs de position}
Pour une série ordonnée \(x_1\le\cdots\le x_N\), \(N\ge1\), la moyenne est \(\bar x=\dfrac1N\sum x_i\); avec des effectifs \(n_j\), elle vaut \(\sum n_jx_j/\sum n_j\). La médiane est la valeur centrale si \(N\) est impair, la moyenne des deux centrales si \(N\) est pair.
\end{frame}
\begin{frame}{Indicateurs de position}
\(Q_1\) est la valeur de rang \(\lceil N/4\rceil\), \(Q_3\) celle de rang \(\lceil3N/4\rceil\). Le décile \(D_k\) est la valeur de rang \(\lceil kN/10\rceil\), \(k=1,\ldots,9\). \(\lceil a\rceil\) désigne le plus petit entier supérieur ou égal à \(a\). Ces conventions sont précisées car elles peuvent différer selon les logiciels.
\end{frame}
\begin{frame}{Dispersion et représentation}
L’étendue vaut maximum moins minimum; l’écart interquartile vaut \(Q_3-Q_1\). La variance descriptive vaut \(V=\dfrac1N\sum(x_i-\bar x)^2\); l’écart type est \(\sigma=\sqrt V\). Une boîte à moustaches représente minimum, \(Q_1\), médiane, \(Q_3\), maximum, selon la convention utilisée ici.
\end{frame}
\begin{frame}{Dispersion et représentation}
\(V=\dfrac1N\sum x_i^2-\bar x^2\). La variance est positive ou nulle, et vaut zéro exactement lorsque toutes les valeurs sont égales.
\end{frame}
\begin{frame}{Dispersion et représentation}
Développer les carrés et utiliser \(\sum x_i=N\bar x\). La première définition est une moyenne de carrés positifs ou nuls; elle est nulle exactement si chaque carré est nul.
\end{frame}
\begin{frame}{Dispersion et représentation}
Série \(4,6,6,8,10,10,12,14\). \(N=8\), somme \(70\), moyenne \(8{,}75\), médiane \(9\); rangs \(2\) et \(6\), donc \(Q_1=6,Q_3=10\). Étendue \(10\), écart interquartile \(4\). Somme des carrés \(692\); variance \(86{,}5-(8{,}75)^2=9{,}9375\); écart type \(\sqrt{9{,}9375}\simeq3{,}152\) minutes. La variance s’exprime en minutes carrées.
\end{frame}
\end{document}
